\section{Introduction}\label{sec:introduction}

Milnor $K$-theory records multiplication in a field together with relations
arising from addition. A reconstruction problem asks whether this
information determines the field, and whether isomorphisms of the
invariant come from field isomorphisms. We study this question in
characteristic $p$ using only the first two Milnor $K$-groups modulo $p$.

Fix a prime $p$ and write $\Lambda=\Fp$. For a field $F$ of
characteristic $p$, set
\[
 V_F=F^\times/(F^\times)^p=\KM_1(F)/p.
\]
We write $[f]$ for the class of $f\in F^\times$ and use additive
notation in $V_F$: thus $[f]+[g]=[fg]$, $r[f]=[f^r]$ for
$r\in\Z$, and $[1]=0$. Define the Steinberg relation subspace
\[
 R_F=\Span_{\Lambda}
       \{[f]\otimes[1-f]:f\in F\setminus\{0,1\}\}
       \subset V_F\otimes_{\Lambda}V_F.
\]
The quotient $W_F=(V_F\otimes_{\Lambda}V_F)/R_F$ is
$\KM_2(F)/p$ by the tensor presentation of Milnor $K$-theory
\cite[Section~1]{Milnor1970}. We denote the image of
$[f]\otimes[g]$ in $W_F$ by $\{f,g\}_F$.

\noindent\textit{Remark (notation).} $F$ denotes the characteristic-$p$ field to which the local field-theoretic definitions or arguments are applied. In Lemma~\ref{lem:p-degree}, Lemma~\ref{lem:zero-symbol}, Lemma~\ref{lem:affine-symbol}, Lemma~\ref{lem:perfect-core}: $F/\kappa$ is finitely generated over an algebraically closed field $\kappa$ of characteristic $p$; no positive lower bound on $\trdeg(F/\kappa)$ is imposed. In Lemma~\ref{lem:p-basis}: $F$ is any field of characteristic $p$; finiteness of $[F:F^p]$ is assumed only for the assertions explicitly conditional on finite degree. In Proposition~\ref{prop:intersection}: $[F:F^p]<\infty$. In Lemma~\ref{lem:scalar-line}: $F\ne F^p$; finite power degree and an algebraically closed base are not assumed. $V_F=F^\times/(F^\times)^p$ has additive operations $[f]+[g]=[fg]$ and $r[f]=[f^r]$, with zero $[1]$, while its underlying set is $\PP_{F^p}(F)$ via $[f]\mapsto F^pf$, whose projective lines come from two-dimensional $F^p$-subspaces of $F$. $g$ is a local nonzero field element; in the class and Milnor-symbol notation it ranges independently over $F^\times$. In Lemma~\ref{lem:perfect-core}: For a given integer $m\ge1$, $g\in F$ is a hypothetical element satisfying $f=g^{p^m}$; it is nonzero because $f\notin\kappa$, and its choice may depend on $m$. For $f\in F^\times$, $[f]$ is its class modulo $(F^\times)^p$ and, under the identification of the underlying set of $V_F$ with $\PP_{F^p}(F)$, represents the one-dimensional $F^p$-subspace $F^p f$; $[1]$ is the zero of the additive group $V_F$. $f$ is a local nonzero element of the field whose multiplicative class is being used. In Lemma~\ref{lem:zero-symbol}: In the logarithmic-derivative pairing, $f$ ranges over $F^\times$; when checking a Steinberg pair $([f],[1-f])$, it ranges over $F\setminus\{0,1\}$. In Theorem~\ref{thm:main}: In the induced-map formula and the comparison of semilinear maps $w\mapsto S(fw)$ and $w\mapsto S(f)S(w)$, $f\in K^\times$. In Lemma~\ref{lem:perfect-core}: In the proof, $f\in F\setminus\kappa$ is chosen, hence is nonzero and transcendental over $\kappa$, and is completed to a transcendence basis. The letter $r$ is local notation; in the convention $r[f]=[f^r]$ it is an integer, whose action on $V_F$ depends only on its residue in $\Fp$. In Lemma~\ref{lem:p-basis}: In the definition of a $p$-independent list $s_1,\ldots,s_r$, $r$ is the finite length of the list and $[C(s_1,\ldots,s_r):C]=p^r$. In Proposition~\ref{prop:intersection}: In the uniqueness argument, $r$ is a candidate integer exponent with $1\le r<p$, satisfying $U^r=h+jX^r$ with $h,j\in C$; the analogous argument applies to $V$. In Lemma~\ref{lem:scalar-line}: $r\in\Fp^\times$ is the scalar multiplying $V_F$; its powers in the proof use its unique integer representative $1\le r<p$. In Proposition~\ref{prop:lines}: In the inverse-map argument, $r\in\Fp^\times$ is chosen so that $r\Psi^{-1}$ maps projective lines into projective lines. In Theorem~\ref{thm:main}: In the injectivity argument, $r\in\Fp^\times$ is the scalar in $\beta_*=r\alpha_*$. $V\in C(Y)^\times$ satisfies $[V]\notin\Fp[Y]$, equivalently $V\notin C^\times Y^i$ for every integer $0\le i<p$. In Proposition~\ref{prop:lines}: With $H=K^p$ and $C=L^p$, $V$ is a representative in $L^\times$ of $\Theta([y+b])$, where $b\in H^\times$ and $[x],[y]$ is the chosen edge. It is determined only up to multiplication by $C^\times$.


Let $\Isom_{\mathrm M}(V_K,V_L)$ be the set of
$\Lambda$-linear isomorphisms $\Theta:V_K\to V_L$ satisfying
\[
                 (\Theta\otimes\Theta)(R_K)=R_L.
\]
We call these isomorphisms \emph{compatible}.
Equivalently, $\Theta$ extends to an isomorphism of the degree-one and
degree-two groups respecting their product pairing. Scalar multiplication
by $\Lambda^\times$ acts on this set. For specified subfields
$k\subset K$ and $l\subset L$, write $\Isom(K,k;L,l)$ for the
field isomorphisms $\alpha:K\to L$ with $\alpha(k)=l$.
Such an isomorphism induces $\alpha_*([f])=[\alpha(f)]$ and
carries $R_K$ onto $R_L$.

\noindent\textit{Remark (notation).} $k$ denotes the specified source subfield $k\subset K$ in the base-field notation and a basis index in the notation $s_k$. In Theorem~\ref{thm:main}, Proposition~\ref{prop:lines}: $k$ is algebraically closed of characteristic $p$, and $K/k$ is finitely generated with $\trdeg(K/k)\ge2$. In Lemma~\ref{lem:p-basis}: $k$ in $s_k$ is an integer basis index with $1\le k\le e$; the intersection assertion is for $j\ne k$. $l\subset L$ is the specified target subfield in $\Isom(K,k;L,l)$. In Theorem~\ref{thm:main}, Proposition~\ref{prop:lines}: $l$ is algebraically closed of characteristic $p$, and $L/l$ is finitely generated with $\trdeg(L/l)\ge2$. $V_K=K^\times/(K^\times)^p$ has its additive $\Lambda$-vector-space structure $[f]+[g]=[fg]$, and its underlying set is $\PP_{K^p}(K)$ via $[f]\mapsto K^pf$, so its projective lines are the point sets of two-dimensional $K^p$-subspaces of $K$. $V_L=L^\times/(L^\times)^p$ has its additive $\Lambda$-vector-space structure $[f]+[g]=[fg]$, and its underlying set is $\PP_{L^p}(L)$ via $[f]\mapsto L^pf$, so its projective lines are the point sets of two-dimensional $L^p$-subspaces of $L$. $\alpha:K\to L$ is a field isomorphism. In Theorem~\ref{thm:main}: As an argument of the canonical map, $\alpha\in\Isom(K,k;L,l)$, so $\alpha(k)=l$.


\begin{theorem}\label{thm:main}
Let $K/k$ and $L/l$ be finitely generated extensions of algebraically
closed fields of characteristic $p$, with
\[
              \trdeg(K/k)\ge2,\qquad \trdeg(L/l)\ge2.
\]
The canonical map
\[
 \Isom(K,k;L,l)\longrightarrow
 \Isom_{\mathrm M}(V_K,V_L)/\Lambda^\times,
 \qquad \alpha\longmapsto[\alpha_*],
\]
is bijective.
\end{theorem}

The invariant therefore determines the field and its isomorphisms, with
exactly the scalar identification in the displayed target. The argument
also recovers the named algebraically closed base from the field itself.

\subsection{Context and method}

Bogomolov and Tschinkel developed a reconstruction method that recovers
projective lines from multiplication and algebraic dependence, and applied
it to Milnor $K$-theory in characteristic zero
\cite{BogomolovTschinkel2009}. Cadoret and Pirutka extended reconstruction
from the multiplicative group modulo constants and algebraic dependence
to finitely generated regular extensions of perfect fields
\cite[Theorem~4]{CadoretPirutka2021}. For algebraically closed or finite
base fields, they also obtained reconstruction from suitable quotients of
the Milnor $K$-ring \cite[Corollary~10]{CadoretPirutka2021}.

Reconstruction from other reduced forms of $K$-theory has also been studied.
For a prime $\ell$ different from the characteristic, Topaz recovers
perfect closures from mod-$\ell$ data together with rational subgroups
for function fields over algebraically closed bases in transcendence
degree at least five \cite[Theorem~B]{Topaz2016}.
His work on rational Milnor $K$-theory reconstructs perfect closures for
fields of absolute transcendence degree at least five
\cite[Main Theorem]{Topaz2023}. Theorem~\ref{thm:main} concerns reduction
modulo the characteristic and starts directly from $V_F$ and $R_F$;
neither rational subgroups nor algebraic-dependence data are supplied.

The geometric framework of our proof has close predecessors. A differential
approach to line recovery appears in an argument of Rovinsky recorded by
Bogomolov and Tschinkel
\cite[Proposition~9]{BogomolovTschinkel2012}. In positive
characteristic, Cadoret and Pirutka recover projective lines over the
subfield of $p$th powers after a common power correction, using a
differential intersection calculation
\cite[Section~1.3, Lemma~23, and Proposition~27]{CadoretPirutka2021}.
Here the intersection calculation involves only degree-$p$ extensions
of the subfield of $p$th powers.
The resulting intersection statement, Proposition~\ref{prop:intersection},
and the passage from vanishing symbols to these subextensions are given
complete elementary proofs below. These are the two structural ingredients
that allow the degree-two mod-$p$ relations to supply the required
projective geometry.

\subsection{Outline of the proof}

Put $C_F=F^p$. Since $(F^\times)^p=C_F^\times$, the underlying
set of $V_F$ is the projective point set $\PP_{C_F}(F)$.
Its point $[1]$ is the zero of the group $V_F$; its projective lines
are the sets of one-dimensional subspaces in two-dimensional
$C_F$-subspaces of $F$. The field degree is finite:
\[
                         [F:C_F]=p^{\trdeg(F/\kappa)}
\]
for $F/\kappa$ as in the theorem. This is a projective geometry over
$C_F$, not over the original constant field $\kappa$.

\noindent\textit{Remark (notation).} $C$ is the scalar field specified in the local field or vector-space setting. In Lemma~\ref{lem:p-degree}, Lemma~\ref{lem:p-basis}, Lemma~\ref{lem:zero-symbol}, Lemma~\ref{lem:affine-symbol}, Proposition~\ref{prop:intersection}, Lemma~\ref{lem:scalar-line}: $C=F^p$. In Theorem~\ref{thm:main}, Proposition~\ref{prop:lines}: $C=L^p$, the power subfield of the target field $L$. In Proposition~\ref{prop:projective-lift}: $C$ is an arbitrary field over which $E'$ is a finite-dimensional vector space of dimension at least three; no power-subfield interpretation is imposed.


The first step constrains pairs with vanishing symbol.
If $s\notin C_F$ and $\{s,t\}_F=0$, then $t\in C_F(s)$.
To prove this implication, consider a $p$-independent pair $s,t$, meaning
that $[C_F(s,t):C_F]=p^2$. Extend it to a monomial basis and take the
determinant of two logarithmic derivatives. This
pairing kills Steinberg relations and is nonzero on the independent
pair. Section~\ref{sec:symbols} supplies the details.

\noindent\textit{Remark (notation).} $s$ denotes a local field element, with its field and restrictions specified by the adjunction, pair, or projective point in which it occurs. In Lemma~\ref{lem:p-basis}: $s\in F\setminus M$, where $C\subseteq M\subseteq F$, is the element in the simple adjunction $M(s)/M$. In Lemma~\ref{lem:zero-symbol}: $s\in F^\times$, and the implication has antecedent $s\notin C$. In the contrapositive proof, $t\notin C(s)$, so $s,t$ is a $p$-independent pair used to choose $D_s,D_t$. In Proposition~\ref{prop:lines}: $s\in K^\times$ represents a source projective point; multiplication by $s^{-1}$ takes the line joining $[s],[t]$ to a line through $[1]$. $t$ denotes an element of the field specified by its local pair, affine expression, or coordinate chart. In Lemma~\ref{lem:zero-symbol}: $t\in F^\times$. In the contrapositive proof, $s\notin C$ and $t\notin C(s)$, so the pair $s,t$ is $p$-independent and is extended to the basis defining $D_s,D_t$. In Lemma~\ref{lem:affine-symbol}: $t\in F^\times$ and $c+dt\ne0$, where $c,d\in C^\times$; no $p$-independence premise is imposed. In Proposition~\ref{prop:lines}: In the line joining $[s],[t]$, $t\in K^\times$ is a representative of a source projective point, and multiplication by $s^{-1}$ translates this line to one through $[1]$. In Proposition~\ref{prop:projective-lift}: $t\in H$ is an arbitrary affine coordinate in the lines $(t,t+a)$ and $(t,at)$; it may be zero.


The fields $C_F(s)$ need not themselves be two-dimensional over $C_F$,
so this implication does not yet recover projective lines. The key
calculation in Section~\ref{sec:intersection} uses intersections of
two multiplicative translates of such subfields. For a $p$-independent
pair $X,Y$, suppose $U\in C_F(X)^\times$ and
$V\in C_F(Y)^\times$ are not scalar multiples of powers of $X,Y$,
respectively. A nonzero element of
\[
             X\,C_F(Y/X)\ \cap\ U\,C_F(V/U)
\]
forces, for a unique $1\le n<p$,
\[
                    U^n\in C_F+C_FX^n,
                    \qquad V^n\in C_F+C_FY^n.
\]
Euler derivations give this common exponent and show that either
inclusion determines it.

\noindent\textit{Remark (notation).} $X$ is a nonzero element of the locally specified ambient field, used through its projective class $[X]$. In Proposition~\ref{prop:intersection}: $X\in F^\times$ and $[C(X,Y):C]=p^2$. In Proposition~\ref{prop:lines}: $X\in L^\times$ represents $\Theta([x])$ for a nonzero vertex $[x]\in V_K$, with $C=L^p$. For an edge $[x],[y]$, $C(X)\ne C(Y)$ and $X,Y$ are $p$-independent over $C$. The representative is determined only up to $C^\times$. $Y$ is a nonzero element of the locally specified ambient field, used through its projective class $[Y]$. In Proposition~\ref{prop:intersection}: $Y\in F^\times$ and $[C(X,Y):C]=p^2$. In Proposition~\ref{prop:lines}: $Y\in L^\times$ represents $\Theta([y])$ for a nonzero vertex $[y]\in V_K$, with $C=L^p$. For an edge $[x],[y]$, $C(X)\ne C(Y)$ and $X,Y$ are $p$-independent over $C$. The representative is determined only up to $C^\times$. $U\in C(X)^\times$ satisfies $[U]\notin\Fp[X]$, equivalently $U\notin C^\times X^i$ for every integer $0\le i<p$. In Proposition~\ref{prop:lines}: With $H=K^p$ and $C=L^p$, $U$ is a representative in $L^\times$ of $\Theta([x+a])$, where $a\in H^\times$ and $[x],[y]$ is the chosen edge. It is determined only up to multiplication by $C^\times$.


Section~\ref{sec:lines} produces the required intersections from
Steinberg relations. Uniqueness propagates the exponent across a
connected graph, giving one scalar $n$ for which $n\Theta$ preserves
every projective line. Applying the argument to the inverse, and using
the fact that no nontrivial scalar preserves projective lines, yields
preservation in both directions. Section~\ref{sec:fields} then proves
the needed projective lifting theorem and uses the multiplicative group
law to turn the normalized semilinear lift into a field isomorphism.
Finally, the identity
$\bigcap_{m\ge1}F^{p^m}=\kappa$ recovers the original base and
completes the theorem.

\noindent\textit{Remark (notation).} $n$ denotes a nonzero power correction modulo $p$, represented by its unique integer in $\{1,\ldots,p-1\}$ when used as an exponent. In Proposition~\ref{prop:intersection}: $n$ is the unique integer $1\le n<p$ such that $U^n\in C+CX^n$ and $V^n\in C+CY^n$; either inclusion determines it uniquely, and both coefficients in each inclusion are nonzero. In Proposition~\ref{prop:lines}, Theorem~\ref{thm:main}: $n\in\Fp^\times$ is the unique scalar such that $\Psi=n\Theta$ sends every projective line over $K^p$ onto a projective line over $L^p$. In the local power calculations its integer representative is used; the scalar is common to all vertices and all affine coefficients.

